\chapter{项目二：实现一个 LRU 缓存}

\section{项目目标}

LRU 缓存是面试题，也是工程中常见结构。需求是：缓存最多保存 \code{capacity} 个键值对；每次访问一个键，这个键变成最新；插入新键时如果容量满，删除最久未使用的键。

一个可用接口如下：

\begin{lstlisting}[language=C++,caption={LRU 缓存接口}]
template <typename Key, typename Value>
class LruCache {
public:
    explicit LruCache(std::size_t capacity);
    [[nodiscard]] std::optional<Value> get(const Key& key);
    void put(Key key, Value value);
    [[nodiscard]] std::size_t size() const noexcept;
};
\end{lstlisting}

\section{为什么不能只用 unordered\_map}

\code{unordered_map} 能快速按键查找，但不知道谁最久没用。只用 \code{vector} 能维护顺序，但查找是线性的。LRU 需要两个结构组合：

\begin{itemize}
  \item \code{std::list} 维护新旧顺序，表头表示最新。
  \item \code{std::unordered_map} 从 key 定位到链表节点。
\end{itemize}

\section{节点和成员}

\begin{lstlisting}[language=C++,caption={LRU 的内部结构}]
#include <list>
#include <optional>
#include <unordered_map>

template <typename Key, typename Value>
class LruCache {
private:
    struct Entry {
        Key key;
        Value value;
    };

    using EntryList = std::list<Entry>;
    using EntryIterator = typename EntryList::iterator;

public:
    explicit LruCache(std::size_t capacity)
        : capacity_{capacity}
    {
    }

private:
    std::size_t capacity_ = 0;
    EntryList entries_;
    std::unordered_map<Key, EntryIterator> index_;
};
\end{lstlisting}

链表迭代器在移动节点时保持有效，这是选择 \code{std::list} 的关键原因。

\section{访问时移动到表头}

\begin{lstlisting}[language=C++,caption={get 操作}]
template <typename Key, typename Value>
std::optional<Value> LruCache<Key, Value>::get(const Key& key)
{
    const auto found = this->index_.find(key);
    if (found == this->index_.end()) {
        return std::nullopt;
    }

    this->entries_.splice(this->entries_.begin(), this->entries_, found->second);
    return found->second->value;
}
\end{lstlisting}

\code{splice} 不复制节点，只把已有节点移动到新位置。移动后迭代器仍指向同一个节点。

\section{插入和淘汰}

\begin{lstlisting}[language=C++,caption={put 操作}]
template <typename Key, typename Value>
void LruCache<Key, Value>::put(Key key, Value value)
{
    if (this->capacity_ == 0) {
        return;
    }

    const auto found = this->index_.find(key);
    if (found != this->index_.end()) {
        found->second->value = std::move(value);
        this->entries_.splice(this->entries_.begin(), this->entries_, found->second);
        return;
    }

    this->entries_.push_front(Entry{.key = std::move(key), .value = std::move(value)});
    this->index_[this->entries_.front().key] = this->entries_.begin();

    if (this->entries_.size() > this->capacity_) {
        const Key evicted_key = this->entries_.back().key;
        this->index_.erase(evicted_key);
        this->entries_.pop_back();
    }
}
\end{lstlisting}

这个版本为清晰牺牲了一点细节：淘汰时复制了 \code{Key}。如果 key 很大，可以先移动或用节点引用优化，但要小心 \code{pop_back} 后引用失效。教材第一版先让所有权和顺序正确。

\section{不变量}

LRU 的核心不变量：

\begin{itemize}
  \item \code{entries_.size() == index_.size()}。
  \item \code{index_} 中每个迭代器都指向 \code{entries_} 中的节点。
  \item 链表表头是最新访问项，表尾是最久未使用项。
  \item 容量为 \code{0} 时，缓存永远为空。
\end{itemize}

测试必须围绕这些不变量设计。

\section{测试用例}

\begin{lstlisting}[language=C++,caption={LRU 行为测试思路}]
LruCache<int, std::string> cache{2};
cache.put(1, "one");
cache.put(2, "two");
assert(cache.get(1) == "one"); // 1 变成最新
cache.put(3, "three");         // 淘汰 2
assert(!cache.get(2).has_value());
assert(cache.get(1) == "one");
assert(cache.get(3) == "three");
\end{lstlisting}

\begin{warningbox}
LRU 最容易错在“查找表和链表不同步”。每次插入、访问、淘汰后，都要能解释两者如何保持一致。
\end{warningbox}

\section{从样例手推状态变化}

先不用代码，手推一个容量为 \code{2} 的缓存。链表从左到右表示“最新到最旧”。

\begin{lstlisting}[caption={LRU 状态推演}]
put(1, one)      [1]
put(2, two)      [2, 1]
get(1) -> one    [1, 2]
put(3, three)    [3, 1]  淘汰 2
get(2) -> miss   [3, 1]
put(1, uno)      [1, 3]  更新 1 并移动到最新
\end{lstlisting}

这段推演给出实现规则：新增节点放表头；命中节点移动到表头；容量超过时删除表尾；更新已有键时不能增加大小。任何实现只要违背其中一条，都会在样例里暴露。

\section{put 的细节拆解}

\code{put} 有三条路径。第一，容量为零，任何插入都无效。第二，键已存在，只更新值并移动节点。第三，键不存在，插入新节点，必要时淘汰旧节点。把这三条路径分清，代码就不会乱。

为了减少淘汰时复制 key，可以先读取表尾 key，再擦除映射，最后弹出节点。注意顺序不能反：

\begin{lstlisting}[language=C++,caption={淘汰顺序必须保证 key 仍可读}]
const Key& evicted_key = this->entries_.back().key;
this->index_.erase(evicted_key);
this->entries_.pop_back();
\end{lstlisting}

这里在 \code{pop_back} 之前使用引用是安全的；\code{pop_back} 之后引用立刻失效，所以不能再访问。教材前面的复制版本更保守，适合先学正确性；这个版本展示如何在清楚生命周期的前提下减少复制。

\section{完整类实现}

把接口、成员和操作合在一起：

\begin{lstlisting}[language=C++,caption={可用的 LRU 缓存实现}]
template <typename Key, typename Value>
class LruCache {
private:
    struct Entry {
        Key key;
        Value value;
    };

    using EntryList = std::list<Entry>;
    using EntryIterator = typename EntryList::iterator;

public:
    explicit LruCache(std::size_t capacity)
        : capacity_{capacity}
    {
    }

    [[nodiscard]] std::optional<Value> get(const Key& key)
    {
        const auto found = this->index_.find(key);
        if (found == this->index_.end()) {
            return std::nullopt;
        }
        this->entries_.splice(this->entries_.begin(), this->entries_, found->second);
        return found->second->value;
    }

    void put(Key key, Value value)
    {
        if (this->capacity_ == 0) {
            return;
        }

        const auto found = this->index_.find(key);
        if (found != this->index_.end()) {
            found->second->value = std::move(value);
            this->entries_.splice(this->entries_.begin(), this->entries_, found->second);
            return;
        }

        this->entries_.push_front(Entry{std::move(key), std::move(value)});
        this->index_[this->entries_.front().key] = this->entries_.begin();
        if (this->entries_.size() > this->capacity_) {
            const Key& evicted_key = this->entries_.back().key;
            this->index_.erase(evicted_key);
            this->entries_.pop_back();
        }
    }

    [[nodiscard]] std::size_t size() const noexcept
    {
        return this->entries_.size();
    }

private:
    std::size_t capacity_ = 0;
    EntryList entries_;
    std::unordered_map<Key, EntryIterator> index_;
};
\end{lstlisting}

\section{为什么返回 optional<Value> 而不是引用}

\code{get} 返回 \code{std::optional<Value>} 会复制值。对于大对象，这可能有成本。为什么教材第一版仍然这么写？因为它简单、安全，调用者拿到的是独立值，不会因为后续缓存淘汰变成悬空引用。

工程里可以设计两个接口：

\begin{lstlisting}[language=C++,caption={按需提供引用接口}]
[[nodiscard]] Value* find(const Key& key);
[[nodiscard]] const Value* find(const Key& key) const;
\end{lstlisting}

指针返回 \code{nullptr} 表示未命中，命中时不复制。但文档必须说明：返回指针只在缓存下一次可能修改前有效。接口越高效，生命周期要求通常越严格。

\section{测试不变量}

除了行为测试，还可以在测试版本里暴露一个检查函数，验证映射和链表同步：

\begin{lstlisting}[language=C++,caption={测试辅助不变量检查}]
[[nodiscard]] bool invariant_holds() const
{
    if (this->entries_.size() != this->index_.size()) {
        return false;
    }
    for (auto it = this->entries_.begin(); it != this->entries_.end(); ++it) {
        const auto found = this->index_.find(it->key);
        if (found == this->index_.end() || found->second != it) {
            return false;
        }
    }
    return true;
}
\end{lstlisting}

生产代码不一定要公开这个函数，但测试思路必须围绕它：每次 \code{put}、\code{get}、淘汰和更新后，链表与哈希表都应该描述同一组条目。

\section{复杂度和替代设计}

\code{get} 和 \code{put} 平均是 \complexity{O(1)}，前提是哈希表表现正常，链表移动是常数时间。替代方案包括：

\begin{itemize}
  \item 只用 \code{vector}：顺序清楚，但查找和移动都是线性。
  \item 只用 \code{unordered_map}：查找快，但不知道谁最旧。
  \item \code{unordered_map} 加自定义双向链表：可控性更高，但实现和测试成本更高。
\end{itemize}

标准库组合已经能满足大多数场景。只有在性能测量证明 \code{std::list} 节点分配成为瓶颈时，才值得考虑自定义存储。

\begin{exercisebox}
给 LRU 增加 \code{contains} 和 \code{erase}。要求 \code{contains} 不改变新旧顺序，\code{erase} 同时删除链表节点和哈希表项。写测试覆盖：删除不存在的 key、删除最新 key、删除最旧 key、删除后再次插入。
\end{exercisebox}
